\documentclass[10pt,a4paper]{amsart}
\usepackage[margin=19mm]{geometry}
\usepackage{amsmath,amssymb,mathtools}
\usepackage[scr=rsfs,cal=euler]{mathalfa}
\usepackage{xfrac}
\usepackage[unicode,hidelinks,bookmarks=false]{hyperref}
\newcommand{\ie}{i.e.\ }
\newcommand{\cf}{cf.\ }
\newcommand{\resp}{respectively\ }
\newcommand{\Subsection}{\subsection}
\providecommand{\nobreakdash}{\nobreak\hbox{-}}
\setlength{\emergencystretch}{2em}
\multlinegap0pt
\usepackage{amssymb}
\usepackage{xcolor}
\newtheorem{theorem}{Theorem}
\DeclareMathOperator{\Var}{Var}
\title{Weak Convergence and Gaussian Limits for a General-Dimensional Origin-Invariant Cram\'er--von Mises Statistic}
\author{Marco Mandap (Bulacan State University, City of Malolos, Philippines)}
\date{Source: arXiv:2605.19164v3 (CC BY 4.0)}
\begin{document}
\maketitle

\noindent\emph{Source and licence.} Weak Convergence and Gaussian Limits for a General-Dimensional Origin-Invariant Cram\'er--von Mises Statistic. Version arXiv:2605.19164v3 (CC BY 4.0), distributed by arXiv under the Creative Commons Attribution 4.0 International licence, http://creativecommons.org/licenses/by/4.0/, as recorded in the arXiv licence metadata for this exact version and shown on the abstract page https://arxiv.org/abs/2605.19164v3. Full licence notice: \texttt{licenses/arxiv-2605.19164-CC-BY-4.0.txt}.\par\smallskip

\noindent\textit{Editorial scope.} Counted unit: the article's theorem thm:op-self-adjoint (source0.tex lines 876--878). The article's complete original proof of it is delivered with it (source0.tex lines 880--895, one \texttt{proof} environment of 16 lines). No mathematics is written, repaired or re-derived: the statement and the proof are the article's own lines, in the article's own notation, and no retained line is substituted or re-flowed.  No further source-local context is needed: every cross-reference of every retained line resolves inside the two delivered spans.  Nothing was omitted from any retained span, so the package carries no omission record.  No retained line contains a citation, so the wrapper carries no bibliography and no bibliography entry.  No retained line contains a figure environment, an image-inclusion command or drawing code, so no image is delivered and the package declares no asset dependency.  Editorial formatting only, in addition: an AMS \texttt{amsart} preamble that restates the article's own declarations and macros used by the retained lines, namely the environments \texttt{theorem} on separate counters, together with the standard packages those lines need.  The retained environments print in this document as Theorem 1 in order of appearance, whereas the article prints the same items with the numbering of its own full document.  The complete native source of the article as distributed in the arXiv e-print for this exact version is bundled unchanged as \texttt{sources/200897/source0.tex}.
\begin{theorem}[Covariance-operator properties]\label{thm:op-self-adjoint}
The operator $T_d$ is bounded, compact, self-adjoint, positive, and trace class.
\end{theorem}

\begin{proof}
Since $K_d$ is bounded and $\nu_d$ is finite, $K_d\in L^2(\nu_d\otimes\nu_d)$; hence $T_d$ is Hilbert--Schmidt and therefore bounded and compact. Symmetry of $K_d$ and Fubini's theorem give self-adjointness. For $f\in L^2(I_d,\nu_d)$,
\[
  \langle T_df,f\rangle
  =\Var_{P_0}\!\left(
    \int_{I_d}f(j)\bigl(\mathbf1_{A_j}(X)-m(j)\bigr)
    \,d\nu_d(j)
  \right)\ge0.
\]
Thus $T_d$ is positive. Finally, the feature map
\[
  X\longmapsto
  \bigl(j\mapsto\mathbf1_{A_j}(X)-m(j)\bigr)
\]
has finite mean squared $L^2(\nu_d)$ norm, so its covariance operator is trace class.
\end{proof}

\end{document}
